% diagrams/firstwords/jar.tex -- el tarro de miel que Andrés le regala a
% Lucía: la tapa de tela con su cordel, la miel y la etiqueta, con una
% abeja que mira por una lupa.
\begin{tikzpicture}[dibujofw]
  % el tarro y la miel
  \filldraw[fill=white, rounded corners=4mm] (-2.6,0) rectangle (2.6,5.6);
  \draw[rounded corners=3mm] (-2.35,0.25) -- (-2.35,4.4) -- (2.35,4.4) -- (2.35,0.25) -- cycle;
  % la tapa de tela, con el cordel
  \filldraw[fill=white] (-2.4,5.2) -- (-2.9,5.9) .. controls (-2.2,6.6) and (2.2,6.6) .. (2.9,5.9) -- (2.4,5.2)
    .. controls (0.8,5.5) and (-0.8,5.0) .. (-2.4,5.2) -- cycle;
  \filldraw[fill=white] (-2.1,5.4) -- (2.1,5.4) -- (2.1,5.65) -- (-2.1,5.65) -- cycle;
  \draw (1.2,5.5) .. controls (1.6,5.0) .. (1.5,4.5); \draw (1.2,5.5) .. controls (1.9,5.2) .. (2.0,4.7);
  % la etiqueta, con la abeja y la lupa
  \filldraw[fill=white] (-1.9,0.9) rectangle (1.9,3.7);
  \begin{scope}[shift={(-0.6,2.3)}, scale=0.62] \abejaFW \end{scope}
  \filldraw[fill=white] (1.05,2.7) circle (0.45); \draw (1.05,2.7) circle (0.28);
  \filldraw[fill=white] (0.75,2.35) -- (0.45,1.95) -- (0.6,1.85) -- (0.9,2.25) -- cycle;
  % la cuchara de la miel
  \filldraw[fill=white] (3.6,0.4) -- (4.2,4.4) -- (4.4,4.4) -- (3.8,0.4) -- cycle;
  \filldraw[fill=white] (3.7,0.6) ellipse (0.45 and 0.75);
  \foreach \y in {0.25,0.6,0.95} {\draw (3.3,\y) -- (4.1,\y);}
\end{tikzpicture}
